\documentclass{article}
\usepackage{amsmath}
\title{Equations and Alignment}
\author{ }
\date{}

\begin{document}
\maketitle

\section*{1. Writing an equation}
A displayed equation can be written between \verb|\[| and \verb|\]|.

\textbf{LaTeX source code:}
\begin{verbatim}
\[
x+2=5
\]
\end{verbatim}
\textbf{Expected compiled output:}
\[
x+2=5
\]

\section*{2. Aligning multiple steps}
The \verb|align*| environment places several lines of maths in columns.

\textbf{LaTeX source code:}
\begin{verbatim}
\begin{align*}
x + 2 &= 5 \\
x &= 5 - 2 \\
x &= 3
\end{align*}
\end{verbatim}
\textbf{Expected compiled output:}
\[
\begin{aligned}
x + 2 &= 5 \\
x &= 5 - 2 \\
x &= 3
\end{aligned}
\]

\textbf{Explanation:}
\begin{itemize}
\item \verb|align*| starts an unnumbered alignment environment.
\item \verb|&| marks the point where columns should line up.
\item \verb|\\| ends a row and starts the next one.
\item The \verb|amsmath| package is required for \verb|align*|.
\end{itemize}

\section*{3. Practice}
Show the steps for solving \(y+4=9\), aligning each equals sign.
\end{document}
